% Format tugas kuliah berbasis pola umum template "homework/assignment"
% (article + identitas mahasiswa + daftar soal), mirip gaya Overleaf
% Professional Assignment / course problem set.
% Soal disalin dari ppt/Logika_Preposisi.tex (bagian Latihan Soal).
\documentclass[12pt,a4paper]{article}

\usepackage[utf8]{inputenc}
\usepackage[T1]{fontenc}
\usepackage[indonesian]{babel}

\usepackage{geometry}
\geometry{margin=2.5cm}

% Font TeX Live/MiKTeX skalabel (menghindari konflik ekspansi font microtype + CM bitmap)
\usepackage{lmodern}

\usepackage{amsmath,amssymb,amsthm}
\usepackage{enumitem}
\usepackage{tabularx}
\usepackage{microtype}

\usepackage{fancyhdr}
\usepackage{lastpage}

\setlength{\headheight}{14pt}
\addtolength{\topmargin}{-2pt}

\usepackage[hidelinks,breaklinks]{hyperref}

% --- Isian identitas (ubah sesuai kebutuhan) ---
\newcommand{\JudulTugas}{Tugas 1}
\newcommand{\KodeMK}{TIF-xxx}
\newcommand{\NamaMK}{Logika Matematika}
\newcommand{\NamaMahasiswa}{[Nama Lengkap Mahasiswa]}
\newcommand{\NIMMhs}{[NIM]}
\newcommand{\Kelas}{[Kelas / Kelompok]}
\newcommand{\Prodi}{Teknik Informatika}
\newcommand{\Fakultas}{Fakultas Teknologi Informasi}
\newcommand{\Institusi}{Universitas Bale Bandung}
\newcommand{\DosenPengampu}{[Nama Dosen Pengampu]}
\newcommand{\TanggalKumpul}{[DD Bulan YYYY]}

\pagestyle{fancy}
\fancyhf{}
\fancyhead[L]{\small\NamaMahasiswa\ --- \NIMMhs}
\fancyhead[R]{\small\JudulTugas\ --- \NamaMK}
\fancyfoot[C]{\small Halaman \thepage\ dari \pageref*{LastPage}}
\renewcommand{\headrulewidth}{0.4pt}
\renewcommand{\footrulewidth}{0pt}

\setlength{\parindent}{0pt}
\setlength{\parskip}{0.6em}

% Jarak vertikal antar lingkungan soal (di atas/di bawah setiap \begin{soal}...\end{soal})
\newtheoremstyle{soaljarak}%
  {1.25ex plus 0.2ex minus 0.1ex}% ruang di atas
  {1.75ex plus 0.35ex minus 0.1ex}% ruang di bawah
  {\normalfont}% isi
  {0pt}%
  {\bfseries}%
  {.}%
  {0.5em}%
  {}
\theoremstyle{soaljarak}
\newtheorem{soal}{Soal}

\begin{document}

\begin{center}
  {\Large\bfseries \JudulTugas}\\[0.35em]
  {\large \NamaMK\ (\KodeMK)}\\[1em]
\end{center}

\noindent
\begin{tabularx}{\linewidth}{@{}p{4.8cm} @{\hspace{0.35em}:\hspace{0.65em}} >{\raggedright\arraybackslash}X@{}}
  Nama Mahasiswa            & \NamaMahasiswa \\
  Nomor Induk Mahasiswa     & \NIMMhs \\
  Program Studi             & \Prodi \\
  Fakultas                  & \Fakultas \\
  Institusi                 & \Institusi \\
  Kelas / Kelompok          & \Kelas \\
  Dosen Pengampu            & \DosenPengampu \\
  Batas tanggal pengumpulan & \TanggalKumpul \\
\end{tabularx}

\vspace{1em}
\hrule
\vspace{1em}

\section*{Petunjuk Umum}
\begin{itemize}[leftmargin=*]
  \item Jawaban \textbf{dianjurkan dikerjakan dalam format \LaTeX}; penggunaan \textit{Microsoft Word} atau aplikasi pengolah kata sejenis \textbf{dibolehkan}.
  \item Setelah selesai, \textbf{simpan atau ekspor jawaban menjadi berkas PDF}.
  \item \textbf{Unggah berkas PDF} tersebut ke \textbf{Edlink} sesuai tenggat di kontrak kuliah/RPS.
  \item Tugas ini membahas \textbf{logika proposisional}. Gunakan huruf proposisional (\textit{p}, \textit{q}, \textit{r}, \ldots) sesuai soal, dan simbol \(\neg\), \(\wedge\), \(\vee\), \(\rightarrow\), \(\leftrightarrow\) secara konsisten. Nilai kebenaran boleh dilambangkan \(T/F\) atau Benar/Salah; jika memakai \(\tau(\cdot)\), tunjukkan langkah substitusinya.
  \item \textbf{Terjemahan bahasa Indonesia \(\leftrightarrow\) bentuk simbol:} perhatikan kata penghubung (``dan'', ``atau'', ``jika \ldots\ maka \ldots'', ``jika dan hanya jika'', ``tidak'') dan susunan yang ekuivalen (misalnya ``\(p\) jika \(q\)'' vs ``jika \(q\) maka \(p\)'').
  \item \textbf{Tabel kebenaran:} lengkapi kolom per sub-ekspresi secara berurutan (sesuai hirarki operator: negasi, konjungsi, disjungsi, implikasi, biimplikasi). Untuk pertanyaan ``apakah sama \ldots?'', bandingkan kolom hasil akhir kedua rumus atau beri argumen ringkas dari tabel.
  \item \textbf{Substitusi nilai \(\tau\):} gantikan setiap proposisi dengan \(T\) atau \(F\), lalu hitung dari dalam kurung ke luar. Sebutkan aturan \(\wedge\), \(\vee\), \(\neg\), \(\rightarrow\), \(\leftrightarrow\) pada setiap langkah jika diperlukan agar jawaban dapat diikuti.
  \item \textbf{Kurung dan ambiguitas:} tuliskan ungkapan yang setara secara jelas dengan tanda kurung, sesuai maksud yang ingin diprioritaskan (lihat urutan hirarki operator pada materi kuliah).
  \item \textbf{Penyederhanaan dan pembuktian ekuivalensi:} tunjukkan rangkaian \(\equiv\) dengan menyebut \textbf{nama hukum} yang dipakai (misalnya identitas, De Morgan, distributif, absorpsi, \(p\rightarrow q \equiv \neg p \vee q\), dan sebagainya). Hindari melompat beberapa langkah tanpa alasan.
  \item \textbf{Tautologi:} untuk Soal 15, sederhanakan ekspresi hingga \(\equiv T\) atau tunjukkan bahwa hasil akhir bukan tautologi (berilah kontra-contoh lewat satu baris tabel kebenaran bila perlu).
  \item Tulis jawaban rapi per nomor soal dan per butir; gambar atau tabel boleh difotokopi rapi dari buku catatan asalkan terbaca.
\end{itemize}

\section*{Soal}

\begin{soal}[\normalfont Sumber: \textit{Logika Preposisi} (slide kuliah)]
  Gunakan konstanta proposisional berikut:
  \begin{center}
    \begin{tabular}{l c l}
      \textit{p} & = & Bowo kaya raya\\
      \textit{q} & = & Bowo hidup bahagia
    \end{tabular}
  \end{center}
  Selanjutnya, ubah pernyataan-pernyataan berikut ini menjadi bentuk logika:
  \begin{enumerate}
    \item Bowo tidak kaya.
    \item Bowo kaya raya dan hidup bahagia.
    \item Bowo kaya raya atau tidak hidup bahagia.
    \item Jika Bowo kaya raya, maka ia hidup bahagia.
    \item Bowo hidup bahagia jika dan hanya jika ia kaya raya.
  \end{enumerate}
\end{soal}

\begin{soal}
  Misalkan \textit{p}, \textit{q} dan \textit{r} adalah variabel proposisional:
  \begin{center}
    \begin{tabular}{l c l}
      \textit{p} & = & Anda sakit flu\\
      \textit{q} & = & Anda ujian\\
      \textit{r} & = & Anda lulus
    \end{tabular}
  \end{center}
  Ubahlah ekspresi berikut menjadi pernyataan dalam bahasa Indonesia:
  \begin{enumerate}
    \item $p \rightarrow \neg q$
    \item $q \rightarrow \neg r$
    \item $\neg q \rightarrow r$
    \item $(p \wedge q) \rightarrow r$
    \item $(p \rightarrow \neg r) \vee (q \rightarrow \neg r)$
    \item $(p \wedge q) \vee (\neg q \wedge r)$
  \end{enumerate}
\end{soal}

\begin{soal}
  Ubahlah pernyataan-pernyataan berikut menjadi bentuk logika:
  \begin{enumerate}
    \item Jika Bowo berada di Malioboro, maka Dewi juga ada di Malioboro.
    \item Pintu rumah Dewi berwarna merah atau coklat.
    \item Berita itu tidak menyenangkan.
    \item Bowo akan datang jika ia mempunyai kesempatan.
    \item Jika Dewi rajin kuliah, maka ia pasti pandai.
  \end{enumerate}
\end{soal}

\begin{soal}
  Dengan menggunakan tabel kebenaran, jawab pertanyaan-pertanyaan berikut:
  \begin{enumerate}
    \item Apakah nilai kebenaran dari $p \wedge p$?
    \item Apakah nilai kebenaran dari $p \vee p$?
    \item Apakah nilai kebenaran dari $(p \wedge \neg p)$ dan $(p \vee \neg p)$?
    \item Apakah $(p \rightarrow q)$ mempunyai nilai kebenaran yang sama dengan $(q \rightarrow p)$?
    \item Apakah $(p \rightarrow q) \rightarrow r$ mempunyai nilai kebenaran yang sama dengan $p \rightarrow (q \rightarrow r)$?
  \end{enumerate}
\end{soal}

\begin{soal}
  Buatlah tabel kebenaran untuk ekspresi-ekspresi logika berikut:
  \begin{enumerate}
    \item $\neg (\neg p \wedge \neg q)$
    \item $p \wedge (p \vee q)$
    \item $((\neg p \wedge (\neg q \wedge r))\vee (q \wedge r))\vee (p \wedge r)$
    \item $(p \wedge q)\vee (((\neg p \wedge q) \rightarrow p)\wedge \neg q)$
    \item $(p \rightarrow q)\leftrightarrow (\neg q \rightarrow \neg p)$
    \item $p \wedge ((r \vee q) \leftrightarrow \neg r)$
    \item $\neg((p \wedge q)\rightarrow \neg r)\vee p$
  \end{enumerate}
\end{soal}

\begin{soal}
  Ubahlah pernyataan-pernyataan berikut menjadi ekspresi logika berupa proposisi majemuk:
  \begin{enumerate}
    \item Jika tikus itu waspada dan bergerak cepat, maka kucing atau anjing itu tidak mampu menangkapnya.
    \item Jika saya tidak keliru, Dewi sudah diwisuda dan pacarnya atau orangtuanya berada di sampingnya.
    \item Bowo membeli saham dan membeli properti untuk investasinya, atau dia dapat menanamkan uang di deposito bank dan menerima bunga uang.
  \end{enumerate}
\end{soal}

\begin{soal}
  Masukkan tanda kurung biasa ke dalam ekspresi logika berikut ini sehingga tidak terjadi ambiguitas:
  \begin{enumerate}
    \item $p \wedge q \wedge r \rightarrow s$
    \item $p \vee q \vee r \leftrightarrow \neg s$
    \item $\neg p \wedge q \rightarrow \neg r \vee s$
    \item $p \rightarrow q \leftrightarrow \neg r \rightarrow \neg s$
    \item $p \vee q \wedge r \rightarrow p \wedge q \vee r$
  \end{enumerate}
\end{soal}

\begin{soal}
  Jika $\tau(p)=T$ dan $\tau(q)=T$, sedangkan $\tau(r)=F$ dan $\tau(s)=F$, carilah nilai kebenaran dari ekspresi-ekspresi logika berikut:
  \begin{enumerate}
    \item $p \wedge (q \vee r)$
    \item $(p \vee q) \wedge r$
    \item $((p\vee q)\wedge r) \vee \neg((p \vee q)\wedge (q\vee s))$
    \item $(\neg(p \wedge q)\vee \neg r) \vee (((\neg p \wedge q)\vee \neg s) \wedge r)$
    \item $(p \leftrightarrow r) \wedge (\neg q \rightarrow s)$
  \end{enumerate}
\end{soal}

\begin{soal}
  Jika $\tau(p)=F$ dan $\tau(q)=F$, sedangkan $\tau(r)=T$ dan $\tau(s)=T$, carilah nilai kebenaran dari ekspresi-ekspresi logika berikut:
  \begin{enumerate}
    \item $\neg(p \leftrightarrow q)\wedge (\neg r \rightarrow s)$
    \item $(\neg p \wedge (q \vee \neg r)\vee (q \leftrightarrow \neg p)) \leftrightarrow (s \wedge r)$
    \item $(p \vee (q \rightarrow (r \wedge \neg p))) \leftrightarrow \neg (q \vee \neg s)$
  \end{enumerate}
\end{soal}

\begin{soal}
  Buktikan bahwa ekspresi-ekspresi logika berikut ini ekuivalen dengan menggunakan tabel kebenaran:
  \begin{enumerate}
    \item $\neg p \leftrightarrow q \equiv (\neg p \vee q)\wedge (\neg q \vee p)$
    \item $q \leftrightarrow (\neg p \rightarrow q) \equiv T$
    \item $(p \vee \neg q)\rightarrow r \equiv (\neg p \wedge q)\vee r$
    \item $p \rightarrow (q \rightarrow r) \equiv (p \rightarrow r)\rightarrow r$
    \item $p \rightarrow q \equiv \neg (p \wedge \neg q)$
    \item $\neg (\neg(p \wedge q)\vee q) \equiv F$
    \item $((p \wedge (q\rightarrow r))\wedge(p \rightarrow(q \rightarrow \neg r)))\rightarrow p \equiv T$
  \end{enumerate}
\end{soal}

\begin{soal}
  Sederhanakan ekspresi-ekspresi logika berikut ini menjadi bentuk paling sederhana:
  \begin{enumerate}
    \item $p \wedge (\neg p \rightarrow p)$
    \item $\neg (\neg p \wedge (q\vee \neg q))$
    \item $(p \rightarrow q)\rightarrow((p \rightarrow \neg q)\rightarrow \neg p)$
    \item $(p \rightarrow(p \vee \neg r))\wedge \neg p \wedge q$
    \item $(\neg p \wedge \neg (q \rightarrow r))\wedge \neg (p \rightarrow (q \rightarrow \neg r))$
    \item $(\neg p \rightarrow q)\rightarrow ((\neg p \rightarrow \neg q)\rightarrow p)$
    \item $(p \wedge (\neg p \vee q))\vee q \vee (p \wedge (p \vee q))$
    \item $(p \wedge \neg q \wedge p \wedge \neg r)\vee (r \wedge p \wedge \neg r)$
  \end{enumerate}
\end{soal}

\begin{soal}
  Buktikan \textit{absorption laws} berikut ini dengan penyederhanaan:
  \begin{enumerate}
    \item $(p \wedge q)\vee(\neg p \wedge q) \equiv q$
    \item $(p \vee q)\wedge(\neg p \vee q) \equiv q$
  \end{enumerate}
\end{soal}

\begin{soal}
  Sederhanakan ekspresi-ekspresi logika berikut ini menjadi bentuk paling sederhana:
  \begin{enumerate}
    \item $\neg p \rightarrow \neg q$
    \item $(p \rightarrow q) \wedge (q \rightarrow r)$
    \item $(p \rightarrow q)\leftrightarrow ((p \wedge q) \leftrightarrow p)$
  \end{enumerate}
\end{soal}

\begin{soal}
  Buktikan dua ekspresi logika berikut ekuivalen dengan penyederhanaan:
  \begin{enumerate}
    \item $(p \wedge q)\vee (q \wedge r) \equiv q \wedge (p \vee r)$
    \item $\neg(\neg(p \wedge q)\vee p) \equiv T$
    \item $\neg(\neg p \vee \neg(r \vee s))\equiv (p \wedge r)\vee (p \wedge s)$
    \item $p \wedge (\neg p \vee q)\equiv p \wedge q$
    \item $\neg(\neg p \wedge \neg q)\equiv p \vee q$
    \item $(p \wedge q)\vee(p \wedge \neg q)\equiv p$
    \item $p \leftrightarrow q \equiv \neg(\neg(\neg p \vee q)\vee \neg(\neg q \vee p))$
    \item $p \leftrightarrow q \equiv (\neg p \vee q) \wedge (\neg q \vee p)$
    \item $\neg(\neg p \vee \neg(q \vee r))\equiv (p \wedge q) \vee (p \wedge r)$
    \item $(p \vee r)\wedge(q \vee s)\equiv (p \wedge q)\vee(p \wedge s)\vee(q \wedge r)\vee(r \wedge s)$
  \end{enumerate}
\end{soal}

\begin{soal}
  Dengan menggunakan penyederhanaan, periksa apakah ekspresi logika berikut merupakan tautologi:
  \begin{enumerate}
    \item $p \wedge (\neg p \rightarrow p)$
    \item $(p \rightarrow q)\rightarrow((p \rightarrow \neg q)\rightarrow \neg p)$
    \item $((p \rightarrow (q \vee \neg r))\wedge \neg p)\wedge q$
    \item $(\neg p \rightarrow q)\rightarrow((\neg p \rightarrow \neg q)\rightarrow p)$
    \item $((p \wedge (q \rightarrow r))\wedge(p \rightarrow(q \rightarrow \neg r)))\rightarrow p$
    \item $((p \rightarrow q)\wedge(q \rightarrow r))\rightarrow((p \wedge q)\rightarrow r)$
    \item $p \rightarrow (q \rightarrow p)$
    \item $(q \rightarrow p) \rightarrow p$
    \item $\neg \neg p \rightarrow p$
    \item $(\neg p \rightarrow \neg q)\rightarrow(q \rightarrow p)$
    \item $(p \rightarrow (q \rightarrow r))\rightarrow((p \rightarrow q)\rightarrow(p \rightarrow r))$
  \end{enumerate}
\end{soal}

\end{document}
